\documentclass[12pt,a4paper]{article}

\usepackage{fontspec}
\usepackage{polyglossia}
\setdefaultlanguage{ukrainian}
\setmainfont{Times New Roman}
\newfontfamily\cyrillicfonttt{Times New Roman}

\usepackage{amsmath,amssymb}
\usepackage{geometry}
\usepackage{float}
\usepackage{caption}
\usepackage{setspace}
\usepackage{hyperref}
\usepackage{graphicx}

\geometry{left=25mm,right=15mm,top=20mm,bottom=20mm}
\setstretch{1.2}
\captionsetup{justification=centering}

\begin{document}

\begin{titlepage}
    \begin{center}
        \textbf{Київський національний університет імені Тараса Шевченка}\\
        \textbf{Факультет комп'ютерних наук та кібернетики}
    \end{center}

    \vspace{4cm}

    \begin{center}
        \textbf{\Large ЗВІТ}\\
        \vspace{0.5cm}
        \textbf{\large до лабораторної роботи №\,3}\\
        \vspace{0.5cm}
        з дисципліни \textit{«Основи нелінійної динаміки»}\\
        \vspace{0.7cm}
        \textbf{Тема: Перевірка умов існування центра\\для системи з квадратичною правою частиною}
    \end{center}

    \vfill

    \begin{flushright}
        Виконав:\\
        Кіщук Ярослав Ярославович
    \end{flushright}

    \vspace{1.5cm}

    \begin{center}
        Київ --- 2026
    \end{center}
\end{titlepage}

\tableofcontents
\newpage

\section{Постановка задачі}

Розглянемо стаціонарну систему звичайних диференціальних рівнянь першого порядку з квадратичною правою частиною:
\begin{equation}\label{eq:system}
\begin{cases}
\dot{x}=y+b x^2+(2c+\alpha)xy+d y^2,\\
\dot{y}=-x+a x^2+(2b+\beta)xy+c y^2.
\end{cases}
\end{equation}
Лінійна частина має вигляд
\[
\dot{x}=y,\qquad \dot{y}=-x,
\]
тому нульовий стан рівноваги для лінійної системи є центром. Мета роботи --- алгоритмічно перевірити умови теореми 5.2, які гарантують, що нульовий стан рівноваги системи~\eqref{eq:system} залишається центром з урахуванням квадратичних членів.

\section{Теорема 5.2}

Для системи~\eqref{eq:system} нульовий стан рівноваги є центром у таких випадках:
\begin{enumerate}
    \item $a+c=0,\quad b+d=0$.
    \item
    \[
    \frac{\alpha}{\beta}=\frac{b+d}{a+c}=k,\qquad
    ak^3-(3b+\alpha)k^2+(3c+\beta)k-d=0.
    \]
    \item $\alpha=0,\quad \beta=0$.
    \item $a=0,\quad c=0,\quad \beta=0$.
    \item
    \[
    b+d=0,\quad \alpha=0,\quad \beta+5a+5c=0,\quad
    ac+2a^2+d^2=0,\qquad a+c\ne 0.
    \]
    \item Після повороту координат
    \[
    x=x_1\cos\varphi-y_1\sin\varphi,\qquad
    y=x_1\sin\varphi+y_1\cos\varphi,
    \]
    де
    \[
    \tan\varphi=-\frac{a+c}{b+d},\qquad b+d\ne 0,
    \]
    отримуємо нові коефіцієнти $a_1,b_1,c_1,d_1,\beta_1$, для яких
    \[
    a_1+c_1=0,\quad \beta_1=0,\quad
    a_1+5b_1+5d_1=0,\quad
    b_1a_1+2a_1^2+d_1^2=0.
    \]
\end{enumerate}

Отже, алгоритм перевірки полягає у послідовній перевірці шести наборів алгебраїчних рівностей. Якщо виконується хоча б один пункт, то за теоремою нульовий стан рівноваги є центром.

\subsection*{Центр у довільній точці $(x_0,y_0)$}

Теорема 5.2 у цьому записі застосовується до особливої точки в початку координат. Якщо потрібно дослідити центр у точці $M=(x_0,y_0)$, спочатку виконуємо зсув
\[
u=x-x_0,\qquad v=y-y_0.
\]
Після підстановки система переписується у змінних $(u,v)$, і перевірка умов теореми виконується для точки $(u,v)=(0,0)$. Отже, центр у точці $M$ в початкових координатах еквівалентний центру в нулі після афінного перенесення.

\section{Алгоритм перевірки}

\subsection{Прямі умови}

Умови 1, 3, 4 та 5 перевіряються прямою підстановкою коефіцієнтів у відповідні рівності. Для числових коефіцієнтів у Python використано малу похибку \texttt{TOL}, щоб уникнути помилок округлення.

\subsection{Умова 2}

Умова 2 містить відношення
\[
\frac{\alpha}{\beta},\qquad \frac{b+d}{a+c}.
\]
Якщо знаменники не дорівнюють нулю, то число $k$ обчислюється безпосередньо:
\[
k=\frac{\alpha}{\beta}=\frac{b+d}{a+c}.
\]
Після цього перевіряється кубічне рівняння
\[
ak^3-(3b+\alpha)k^2+(3c+\beta)k-d=0.
\]

Щоб алгоритм не ламався у вироджених випадках, у ноутбуку використано еквівалентну форму без ділення:
\begin{equation*}
\alpha=\beta k,\qquad b+d=(a+c)k.
\end{equation*}
Тому перевірка зводиться до пошуку дійсного $k$, яке одночасно задовольняє дві лінійні рівності та кубічне рівняння.

\subsection{Умова 6}

Для перевірки умови 6 виконується поворот координат. Позначимо
\[
R=
\begin{pmatrix}
\cos\varphi & -\sin\varphi\\
\sin\varphi & \cos\varphi
\end{pmatrix}.
\]
Якщо $F(x,y)$ --- права частина системи~\eqref{eq:system}, то в нових координатах $(x_1,y_1)$ векторне поле обчислюється за формулою
\[
G(x_1,y_1)=R^T F\bigl(R(x_1,y_1)\bigr).
\]
Після розкриття дужок воно знову має канонічний вигляд:
\begin{equation*}
\begin{cases}
\dot{x}_1=y_1+b_1x_1^2+(2c_1+\alpha_1)x_1y_1+d_1y_1^2,\\
\dot{y}_1=-x_1+a_1x_1^2+(2b_1+\beta_1)x_1y_1+c_1y_1^2.
\end{cases}
\end{equation*}
Звідси виділяються коефіцієнти $a_1,b_1,c_1,d_1,\beta_1$ і перевіряються рівності з пункту 6.

\section{Програмна реалізація}

Перевірку реалізовано у Jupyter-ноутбуку:
\[
\texttt{theorem\_5\_2\_checks.ipynb}.
\]
На початку ноутбуку задаються коефіцієнти:
\begin{verbatim}
a = 1
b = Rational(1, 5)
c = Rational(-3, 10)
d = Rational(2, 5)
alpha = 0
beta = 0
\end{verbatim}
Після цього викликається функція \texttt{check\_all(coeffs)}, яка повертає таблицю результатів для шести умов теореми.

Фрагмент реалізації:
\begin{verbatim}
checks = check_all(coeffs)
if any(ok for _, ok, _, _ in checks):
    true_conditions = [number for number, ok, _, _ in checks if ok]
    print("Центр. Виконані умови:", true_conditions)
else:
    print("Жодна з умов теореми 5.2 не виконалась.")
\end{verbatim}

Додатково в ноутбуку побудовано фазові портрети:
\begin{itemize}
    \item повний портрет для нелінійної системи~\eqref{eq:system};
    \item портрет лінеаризованої системи $\dot x = y,\ \dot y = -x$ в околі $(0,0)$.
\end{itemize}
Візуалізація виконується через \texttt{streamplot}, що дозволяє порівняти локальну картину біля особливої точки з глобальнішою геометрією траєкторій на обраній області.

\section{Фазові портрети}

На рис.~\ref{fig:phase-portraits} наведено фазові портрети, згенеровані у ноутбуці \texttt{theorem\_5\_2\_checks.ipynb} для тих самих коефіцієнтів системи. Ліворуч показано повне нелінійне поле напрямків, праворуч --- лінеаризовану систему в околі $(0,0)$.

\begin{figure}[H]
    \centering
    \includegraphics[width=0.98\textwidth]{/Users/user/uni/10-sem/OND/lab3/phase_portraits_theorem_5_2.png}
    \caption{Фазові портрети: нелінійна система (ліворуч) і лінеаризована система (праворуч).}
    \label{fig:phase-portraits}
\end{figure}

\section{Приклад перевірки}

Для прикладових коефіцієнтів
\[
a=1,\quad b=\frac15,\quad c=-\frac{3}{10},\quad d=\frac25,\quad
\alpha=0,\quad \beta=0
\]
маємо
\[
\alpha=0,\qquad \beta=0.
\]
Отже, виконується умова 3 теореми 5.2. Тому за теоремою нульовий стан рівноваги системи~\eqref{eq:system} є центром.

\section{Висновок}

У роботі побудовано алгоритм перевірки умов теореми 5.2 для системи з квадратичною правою частиною. Умови 1, 3, 4, 5 перевіряються прямою підстановкою. Умова 2 реалізована через пошук числа $k$ без обов'язкового ділення на потенційно нульові знаменники. Умова 6 перевіряється після ортогонального повороту координат та виділення нових коефіцієнтів квадратичної частини.

Якщо хоча б одна з шести умов виконується, то за теоремою нульовий стан рівноваги є центром. Якщо жодна умова не виконується, це не означає автоматично фокус: у межах теореми 5.2 отримуємо лише факт, що центр не доведено цим критерієм.

\end{document}
